\renewcommand{\arraystretch}{1.2}
\captionsetup{justification=raggedright,singlelinecheck=false}
\begin{table}[H]
\centering
\caption{Analisis Estimasi Biaya Opsional dan Kondisi Pemicu}
\label{tbl:cost_analysis}
\begin{tabular}{|p{3cm}|p{3cm}|p{7cm}|}
\hline
\textbf{Komponen Opsional} &
\textbf{Estimasi Biaya} &
\textbf{Kondisi Pemicu (\textit{Trigger Condition})} \\ \hline

Layanan LLM Cloud (OpenAI GPT--4o-mini) &
Rp 50.000 -- Rp 100.000 \newline 
(\textit{Pay--per--token}) &
\begin{enumerate}[leftmargin=0.5cm]
    \item Beban memori tinggi pada laptop (16\,GB) ketika menjalankan Model Lokal 
    dan Neo4j secara bersamaan sehingga menyebabkan \textit{system crash} atau 
    \textit{Out-of-Memory (OOM)}. 
\end{enumerate} \\ \hline

Cloud GPU (Google Colab Pro) &
\(\sim\) Rp 160.000 / bulan &
\begin{enumerate}[leftmargin=0.5cm]
    \item Proses \textit{embedding}/vektorisasi data melebihi 12 jam pada CPU laptop 
    sehingga menghambat iterasi dan eksplorasi eksperimen.
    \item Diperlukan pengujian model pembanding berukuran besar 
    ($>$ 14B parameter) yang tidak dapat dijalankan pada VRAM laptop.
\end{enumerate} \\ \hline

\textbf{Total Dana Siaga} &
\(\sim\) Rp 250.000 -- Rp 400.000 &
Dana hanya digunakan apabila skenario eksekusi di lingkungan lokal tidak mampu berjalan 
secara stabil, efisien, atau menghasilkan keluaran yang memenuhi standar validitas. 
Jika seluruh eksperimen dapat dijalankan pada lingkungan lokal, 
maka dana tidak direalisasikan. \\ \hline

\end{tabular}
\end{table}